% GENERATED FILE. Edit canonical lessons, then run npm run generate:books.
% canonical-source-hash: fnv1a64:49eb20e9d16fe93c
% canonical-lessons: LA-C402-edisco, LA-C402-perdisco, LA-C402-studeo, LA-C402-meditor, LA-C402-recordor

\chapter{Learning by Heart}
\label{ch:la-learning-by-heart}
\hlchaptermodality{402}

\begin{chapteropening}
\textbf{By the end of this chapter:} \emph{I can say that I learn something by heart or thoroughly, that I study and rehearse it, and that I remember it.}

Complete the last lesson of chapter 402: I can say that I learn something by heart or thoroughly, that I study and rehearse it, and that I remember it.
\end{chapteropening}

\section[ēdiscō, ēdiscere]{ēdiscō, ēdiscere — to learn by heart}

\label{lesson:LA-C402-edisco}



\begin{quote}
Before the new one: say the Latin for an elementary teacher of letters, then the Latin for a schoolfellow.
\end{quote}

\subsection*{You'll want to know: ēdiscō, ēdiscere}
\textbf{ēdiscō, ēdiscere} — \textquotedblleft{}to learn by heart\textquotedblright{}.

\begin{quote}
Puerī carmina poētae ēdiscunt.
\end{quote}

\begin{quote}
\emph{The boys learn the poet's poems by heart.}
\end{quote}

\emph{ēdiscō, ēdiscere, ēdidicī}: \emph{ē} + \emph{discō}, to learn out to the end.

\begin{cousinweb}
\emph{ēdiscō} is built from \emph{ē} \textquotedblleft{}out, to the end\textquotedblright{} and \emph{discō} \textquotedblleft{}to learn\textquotedblright{} (Lewis \& Short): to learn right through.
\end{cousinweb}

\begin{grammarlens}[title={What you've built}]
One more word for telling, explaining and describing at B1.
\end{grammarlens}

\subsection*{Your turn}
\begin{itemize}
\raggedright
  \item \emph{Say it:} \emph{ēdiscō, ēdiscere}
  \item \emph{Say it:} \emph{ēdiscō, ēdiscere}, once more
  \item \emph{Say it:} a sentence of your own with \emph{ēdiscō, ēdiscere}
  \item \emph{Recall:} say the Latin for an elementary teacher of letters, then the Latin for a schoolfellow, then say \emph{ēdiscō, ēdiscere} again
\end{itemize}

\begin{tcolorbox}[breakable,colback=teal!4,colframe=teal!35!black,title={Before you move on}]
What does ēdiscō, ēdiscere mean? (\textquotedblleft{}To learn by heart\textquotedblright{}.) Say it once more.
\end{tcolorbox}
\section[perdiscō, perdiscere]{perdiscō, perdiscere — to learn thoroughly}

\label{lesson:LA-C402-perdisco}



\begin{quote}
Before the new one: say the Latin for a schoolfellow, then the Latin for to learn by heart.
\end{quote}

\subsection*{You'll want to know: perdiscō, perdiscere}
\textbf{perdiscō, perdiscere} — \textquotedblleft{}to learn thoroughly\textquotedblright{}.

\begin{quote}
Artem medicīnae perdidicit.
\end{quote}

\begin{quote}
\emph{He learned the art of medicine thoroughly.}
\end{quote}

\emph{perdiscō, perdiscere, perdidicī}: \emph{per-} through + \emph{discō}.

\begin{cousinweb}
\emph{perdiscō} is built from \emph{per-} \textquotedblleft{}through, thoroughly\textquotedblright{} and \emph{discō} \textquotedblleft{}to learn\textquotedblright{} (Lewis \& Short).
\end{cousinweb}

\begin{grammarlens}[title={What you've built}]
One more word for telling, explaining and describing at B1.
\end{grammarlens}

\subsection*{Your turn}
\begin{itemize}
\raggedright
  \item \emph{Say it:} \emph{perdiscō, perdiscere}
  \item \emph{Say it:} \emph{perdiscō, perdiscere}, once more
  \item \emph{Say it:} a sentence of your own with \emph{perdiscō, perdiscere}
  \item \emph{Recall:} say the Latin for a schoolfellow, then the Latin for to learn by heart, then say \emph{perdiscō, perdiscere} again
\end{itemize}

\begin{tcolorbox}[breakable,colback=teal!4,colframe=teal!35!black,title={Before you move on}]
What does perdiscō, perdiscere mean? (\textquotedblleft{}To learn thoroughly\textquotedblright{}.) Say it once more.
\end{tcolorbox}
\section[studeō, studēre]{studeō, studēre — to be eager (for); to study}

\label{lesson:LA-C402-studeo}



\begin{quote}
Before the new one: say the Latin for to learn by heart, then the Latin for to learn thoroughly.
\end{quote}

\subsection*{You'll want to know: studeō, studēre}
\textbf{studeō, studēre} — \textquotedblleft{}to be eager (for); to study\textquotedblright{}.

\begin{quote}
Litterīs studet.
\end{quote}

\begin{quote}
\emph{He is devoted to literature.}
\end{quote}

\emph{studeō, studēre, studuī} + dative: \emph{litterīs studēre}, to give oneself to literature. \emph{Studium} is its noun. English \emph{student} and \emph{study}.

\begin{cousinweb}
\emph{studeō} is usually traced to a Proto-Indo-European root meaning \textquotedblleft{}to push, to strike\textquotedblright{} (EtymDB): to press on towards a thing. English \emph{student} and \emph{study} come from it.
\end{cousinweb}

\begin{grammarlens}[title={What you've built}]
One more word for telling, explaining and describing at B1.
\end{grammarlens}

\subsection*{Your turn}
\begin{itemize}
\raggedright
  \item \emph{Say it:} \emph{studeō, studēre}
  \item \emph{Say it:} \emph{studeō, studēre}, once more
  \item \emph{Say it:} a sentence of your own with \emph{studeō, studēre}
  \item \emph{Recall:} say the Latin for to learn by heart, then the Latin for to learn thoroughly, then say \emph{studeō, studēre} again
\end{itemize}

\begin{tcolorbox}[breakable,colback=teal!4,colframe=teal!35!black,title={Before you move on}]
What does studeō, studēre mean? (\textquotedblleft{}To be eager (for); to study\textquotedblright{}.) Say it once more.
\end{tcolorbox}
\section[meditor, meditārī]{meditor, meditārī — to think over; to rehearse, to practise}

\label{lesson:LA-C402-meditor}



\begin{quote}
Before the new one: say the Latin for to learn thoroughly, then the Latin for to be eager.
\end{quote}

\subsection*{You'll want to know: meditor, meditārī}
\textbf{meditor, meditārī} — \textquotedblleft{}to think over; to rehearse, to practise\textquotedblright{}.

\begin{quote}
Ōrātiōnem tōtam noctem meditābātur.
\end{quote}

\begin{quote}
\emph{He rehearsed his speech all night.}
\end{quote}

Deponent: \emph{meditor, meditārī, meditātus sum}. To think over, and to practise, like a speaker rehearsing. English \emph{meditate}.

\begin{cousinweb}
\emph{meditor} is usually traced to the same root as \emph{medeor} \textquotedblleft{}to heal\textquotedblright{} (Lewis \& Short compare \emph{medeor} and \emph{remedium}), a root meaning \textquotedblleft{}to take fitting measures\textquotedblright{}. English \emph{meditate} comes from it.
\end{cousinweb}

\begin{grammarlens}[title={What you've built}]
One more word for telling, explaining and describing at B1.
\end{grammarlens}

\subsection*{Your turn}
\begin{itemize}
\raggedright
  \item \emph{Say it:} \emph{meditor, meditārī}
  \item \emph{Say it:} \emph{meditor, meditārī}, once more
  \item \emph{Say it:} a sentence of your own with \emph{meditor, meditārī}
  \item \emph{Recall:} say the Latin for to learn thoroughly, then the Latin for to be eager, then say \emph{meditor, meditārī} again
\end{itemize}

\begin{tcolorbox}[breakable,colback=teal!4,colframe=teal!35!black,title={Before you move on}]
What does meditor, meditārī mean? (\textquotedblleft{}To think over; to rehearse, to practise\textquotedblright{}.) Say it once more.
\end{tcolorbox}
\section[recordor, recordārī]{recordor, recordārī — to call to mind, to remember}

\label{lesson:LA-C402-recordor}



\begin{quote}
Before the new one: say the Latin for to be eager, then the Latin for to think over.
\end{quote}

\subsection*{You'll want to know: recordor, recordārī}
\textbf{recordor, recordārī} — \textquotedblleft{}to call to mind, to remember\textquotedblright{}.

\begin{quote}
Diem illum saepe recordor.
\end{quote}

\begin{quote}
\emph{I often recall that day.}
\end{quote}

Deponent: \emph{recordor, recordārī, recordātus sum}, built on \emph{cor}, heart: to bring back to the heart. English \emph{record} comes from it.

\begin{cousinweb}
\emph{recordor} is built from \emph{re-} \textquotedblleft{}back\textquotedblright{} and \emph{cor, cordis} \textquotedblleft{}heart\textquotedblright{} (EtymDB; Lewis \& Short), the seat of memory for the Romans. English \emph{record} comes from it through French.
\end{cousinweb}

\begin{grammarlens}[title={What you've built}]
One more word for telling, explaining and describing at B1.
\end{grammarlens}

\subsection*{Your turn}
\begin{itemize}
\raggedright
  \item \emph{Say it:} \emph{recordor, recordārī}
  \item \emph{Say it:} \emph{recordor, recordārī}, once more
  \item \emph{Say it:} a sentence of your own with \emph{recordor, recordārī}
  \item \emph{Recall:} say the Latin for to be eager, then the Latin for to think over, then say \emph{recordor, recordārī} again
\end{itemize}

\begin{tcolorbox}[breakable,colback=teal!4,colframe=teal!35!black,title={Before you move on}]
What does recordor, recordārī mean? (\textquotedblleft{}To call to mind, to remember\textquotedblright{}.) Say it once more.
\end{tcolorbox}
